\section{The Question: One Crisis, or Several Histories?}

Four assertions are often compressed into one sentence: Vatican I made the papacy absolute; Vatican II used that papacy to revolutionize Catholic life; Paul VI replaced the Roman Mass; and Catholic belief and religious life then collapsed. Each clause touches something real. Joined without distinctions, however, they become a story too simple to be true and too powerful to be harmless. It can turn dogma into sociology, a council into its reception, a missal into every parish experiment, and grief over lost vocations into a license for schism.

This article asks a harder question. What did the two Vatican Councils actually do? What kind of authority belongs to each act? How did the Missal of Paul VI arise, and which features came from the council, from the postconciliar commission, from papal decisions, from episcopal conferences, or from later custom? What precisely declined, where, and when? What do surveys really establish about lay belief? Finally, by what tested mechanisms may any of these histories be called causes of the others?

\articlethesis{Vatican I, Vatican II, Paul VI's liturgical reform, and the post-1965 Catholic crisis belong to one history but are not one act or one cause. Vatican I defined true doctrine in a formula whose careful limits were obscured by the council's interruption and by maximalist reception. Vatican II restored much of the missing ecclesial setting and taught authoritatively without making a new extraordinary dogmatic definition. Paul VI's Missal was a lawful postconciliar papal reform implementing, but not simply identical with, the council's mandate. Severe Western contraction and weak doctrinal transmission are facts; conciliar and postconciliar choices are among their institutional causes and accelerants, but the evidence does not make either council or either missal a sufficient cause of a worldwide collapse.}

\subsection{What this study affirms before it criticizes}

A sober Catholic history is neither an acquittal brief nor a prosecution. It begins inside the Church's rule of faith.

\begin{itemize}
\item The Roman primacy and the ex cathedra infallibility defined by Vatican I are divinely revealed dogmas. Their truth is not contingent on the political wisdom of Pius IX, the completeness of the 1870 agenda, or later statistical outcomes.
\item Vatican II was a valid ecumenical council approved and promulgated by the Roman Pontiff. “Pastoral” does not mean private advice, and the absence of new extraordinary definitions does not reduce its authentic teaching to opinion.
\item The Roman Missal promulgated by Paul VI is a valid, lawful liturgical book of the Roman Rite. Authority and sacramental validity do not make every design choice, translation, implementation, or local practice beyond historical and prudential criticism.
\item Numerical growth does not prove fidelity, and numerical decline does not disprove doctrine. Christ did not promise every Catholic institution an uninterrupted increase. Yet pastors and historians may not use that truth to evade catastrophic evidence about lost practice, vocations, or transmission.
\item Communion with the Church is not preserved by rejecting the authority one claims to defend. The traditionalist rupture must be judged by the same exact doctrine of primacy that its adherents invoke against modern Rome.
\end{itemize}

These affirmations make criticism more exact, not less. A dogmatic formula can be true and historically compressed. A council can be authoritative and contain prudential judgments. A lawful reform can be pastorally damaging in some dimensions or implemented with destructive haste. A demographic crisis can have both external pressures and internal ecclesial causes.

\newcommand{\crisisapparatusappendix}{%
\section{Objects, Method, and Qualifications}

\subsection{Six objects that must not be collapsed}

\begin{threestable}{Object}{What it is}{What it is not}
Vatican I's definitions & The two promulgated dogmatic constitutions, especially the defined clauses of \emph{Dei Filius} and \emph{Pastor aeternus}. & Every ultramontane slogan, every private view of Pius IX, or a completed treatise on the whole Church.\\
Vatican II's texts & Sixteen papally promulgated conciliar documents, read by genre, words, subject, and stated theological note. & The whole “council-event,” the media image, every expert's agenda, or every later innovation advertised as its spirit.\\
The liturgical mandate & The principles and directives of \emph{Sacrosanctum Concilium}, especially nos.~21--23, 36--40, and 47--58. & The complete text of the 1969 Order of Mass or 1970 Missal, which did not yet exist.\\
The postconciliar reform & A sequence of papal, curial, commission, episcopal-conference, translation, and local decisions. & One instantaneous act, one language, or every practice that appeared after 1965.\\
The Catholic crisis & Specified changes in attendance, sacraments, vocations, religious institutes, schools, affiliation, belief, and institutional trust. & One universal curve or proof that grace deserted the Church.\\
Traditionalist responses & Several distinguishable canonical communities, associations, attitudes, and individual choices. & A single movement with one legal status, or a license to call every attached Catholic schismatic.\\
\end{threestable}

\subsection{How causation will be judged}

“After” is not yet “because of.” But temporal proximity is not nothing: when a large institution changes many of its authoritative signals and then loses members and practice rapidly, causal inquiry is warranted. The proper response is not to ban the question but to specify it.

\begin{fourstable}{Causal term}{Meaning here}{Evidence needed}{Common mistake}
Necessary cause & Without it, the result could not have occurred. & A defensible counterfactual across comparable settings. & Treating one visible event as indispensable merely because it came first.\\
Sufficient cause & It can produce the result without the other proposed factors. & Mechanism plus cases where it operates independently. & Claiming the council or the sexual revolution alone explains everything.\\
Contributing cause & It changed the probability, magnitude, direction, or timing of the result. & Timing, mechanism, comparative variation, and rival explanations. & Demanding monocausality before admitting any responsibility.\\
Mediator or accelerator & It transmitted external pressure through an institution or made an existing trend faster. & Evidence about leadership, rules, symbols, incentives, and implementation. & Saying external secularization excuses internal decisions, or vice versa.\\
Marker & It coincided with a deeper change without materially causing it. & Failure to find a mechanism after serious testing. & Mistaking the most memorable symbol for the operative cause.\\
\end{fourstable}

The causal unit also matters. “Vatican II caused decline” may refer to a promulgated sentence, the expectation of unlimited change, a national implementation, a religious chapter, a seminary pedagogy, a parish liturgy, or the use of papal authority to suppress inherited practice. Those are different propositions. This study accepts a causal claim only at the level for which evidence exists.

\subsection{Source boundary and intended reader}

Conciliar and magisterial claims are checked against official Holy See texts and their exact loci. Liturgical history distinguishes typical editions, decrees, effective dates, translations, participant accounts, and later discipline. Canonical claims apply the 1983 Latin Code and competent Roman acts through 16 July 2026. Counts come from CARA's series using the \emph{Official Catholic Directory} and Vatican statistical yearbooks; belief claims come from named survey instruments whose question wording is retained. Professional historians and sociologists supply reconstruction and rival causal models, not a substitute magisterium.

The reader should therefore expect unequal conclusions. The definition of papal infallibility receives an assent that no demographic inference can receive. An official count can establish how many women religious were reported, but not why they left. A survey can reveal how respondents answered, but not whether the respondent understood a dogmatic formula, possessed supernatural faith, or was culpable for error. The article's causal synthesis is an argument offered for inspection.

The work is deliberately bounded. It is not the later full comparison of the 1962 and postconciliar forms of the Roman Mass. It does not adjudicate an individual's canonical status or tell a person where to worship. Its aim is prior: to construct a history strong enough that later judgments need not be built from slogans.

\judgmentbox{The crisis is too grave for apologetic denial and too complex for a scapegoat. Fidelity requires holding together truths that partisan histories divide: the councils' authority, the reform's lawful character, the possibility of failed prudence, the reality of institutional contraction, the force of external secularization, the agency of Catholic leaders, and the sin of rupturing communion.}
}
